\documentclass[5pt, a4paper]{article}

\usepackage[utf8]{inputenc}
\usepackage[T1]{fontenc}
\usepackage{textcomp}
\usepackage{amsmath, amssymb}
\usepackage{multicol}
\usepackage{proof}
\usepackage{enumitem}
\usepackage{fancyhdr}
\usepackage[margin=0.5in]{geometry}
% figure support
\usepackage{import}
\usepackage{xifthen}
\pdfminorversion=7

\usepackage[most]{tcolorbox}
\usepackage{pdfpages}
\usepackage{transparent}
\newcommand{\incfig}[1]{%
    \def\svgwidth{\columnwidth}
    \import{./figures/}{#1.pdf_tex}
}

\tcbset{
colback=gray!5,
colframe=blue!50!white,
coltitle=black,
arc=0pt,
boxrule=-0.1mm,
fonttitle=\bfseries,
coltitle=black,
left=2.5mm,
leftrule=1mm,
right=3.5mm,
colbacktitle=black!10!white,
toptitle=0.75mm,
bottomtitle=0.5mm,
}
\begin{document}
\begin{multicols}{2}
    \subsection*{Common Patterns when dealing with Limits}
    \begin{enumerate}
        \item Polynomials.
            Solve them as normal. Try to simplify them as much as possible.
        \item Triginometric.
            Try to split the expression into sine/consine, and match the input
            in the numerator and denominator.
        \item L'Hopital.
            Only use when subtituting results in an indeterminate form.
        \item Difference in Square Roots.
            Say you have \[
                \lim_{n \to \infty}\left( \sqrt{n - 1} - n \right) 
            .\] 
            You usually want to multiply and divide the expression by the sum of
            the square root. This would make the expression much easier to
            evaluate.
        \item Expression with exponentiation.
            Take the logarithm of both sides to bring the exponent down.
        \item $1^\infty$ form.
            Usually can be simplified to Euler's number limit definition.
    \end{enumerate}
    \subsection*{General Framework for proving a Convergence for a Monotonic
    Sequence}
    \begin{enumerate}
        \item Come up with some candidate limits where the sequence converge.
        \item Show that the sequence is monotonic using induction
        \item Show that it is bounded using induction.
        \item Apply Theorem and conclude
    \end{enumerate}
    \begin{tcolorbox}[title=Squeeze Theorem]
        Let $a_n,b_n,c_n$ be sequences. If  \[
        b_n \le  a_n \le c_n
        \] 
        for all $n$, and $b_n \to L$, $c_n \to L$, then $a_n \to L$.
    \end{tcolorbox}
    \begin{tcolorbox}[title=Continuous Function Theorem for Sequences]
        Let $\{a_n\}$ be a sequence. If $a_n \to L$ and $f$ is a function that
        is continuous at $L$, then $f(a_n) \to f(L)$.
    \end{tcolorbox}
    \begin{tcolorbox}[title=Monotone Convergence Theorem]
       For a monotonous sequence (those that is non-decreasing/non-increasing)
       that is bounded, the sequence is convergent.
   \end{tcolorbox} 
   \begin{tcolorbox}[title=Commonly Found Limits]
       \begin{enumerate}
           \item $
           lim_{n \to \infty} \frac{\ln n}{n} = 0 \\
           $
        \item $
            lim_{n \to \infty} x^{\frac{1}{n}} = 0 \\
            $
        \item$ 
            lim_{n \to \infty} \left( 1 + \frac{x}{n} \right)^n  = e^x \\
            $
        \item $ 
            lim_{n \to \infty} \sqrt{n}^n  = 1 \\
            $
        \item $ 
            lim_{n \to \infty} x^n  = 0, \quad |x| < 1 \\
            $
        \item $
            lim_{n \to \infty} \frac{x^n}{n!}  = 0 \\
            $
       \end{enumerate}
   \end{tcolorbox}
   \begin{tcolorbox}[title=Geometric Series]
       In general, an infinite geometric series that converges is given by \[
           \sum_{n=1}^{\infty} ar^{n-1} = \frac{a}{1-r}
       .\] 
       A geometric series converges only if $|r| < 1$. Otherwise, it diverges.
   \end{tcolorbox}

   \subsection*{General Framework for Testing Series Convergence}
  \begin{enumerate}
      \item Non-alternating Series
          \begin{enumerate}
              \item n-th Term Divergence Test. 
                  Always start with this. If the result is not zero, then immediately
                  divergent. Otherwise, proceed.
              \item Ratio/Root Test.  
                  Ratio test is good for factorials, while Root is good for
                  exponentials
              \item Comparison Test or Limit Comparison Test. Useful for Expression
                  with an oscilating term.
              \item Alternating Series Test.
          \end{enumerate}
      \item Alternating Series
          \begin{enumerate}
              \item n-th Term Divergence Test. 
              \item Alternating Series Test. 
                \item Ration/Root Test for checking absolute convergence.
          \end{enumerate}
  \end{enumerate} 
   \begin{tcolorbox}[title= n-th Term Convergence Test]
       A series $\sum a_n$ converges if and only if $a_n \to 0$. A sequence diverge if
       $\lim_{n \to \infty} a_n$ does not exist or is different from zero.
   \end{tcolorbox}
   \begin{tcolorbox}[title=The Integral Test (Extra JIC)]
       Let $\{a_n\} $ be a sequence of positive terms. Suppose we model $\{a_n\}
       = f(n)$ that is continuous, positive, decreasing function of $x$ for all
       $x \ge N$. Then, the series $\sum_{n=N} a_n$ and $\int_N a_n$ both
       diverge or both converge.
   \end{tcolorbox}
   \begin{tcolorbox}[title=Comparison Test]
       Let $\sum a_n, \sum b_n, \sum c_n$ be some series with nonnegative terms
       that for some integer $N$, \[
       b_n \le a_n \le c_n
       \] 
       for all $n \ge N$. Then, if 
       \begin{enumerate}
           \item $\sum c_n$ converge, then $\sum a_n$ converge as well.
            \item $\sum b_n$ diverge, then $\sum a_n$ diverge as well.
       \end{enumerate}
   \end{tcolorbox}
   \begin{tcolorbox}[title=Limit Comparison Test]
       Let $a_n > 0$ $b_n >0$ for all $n$. We define $c$ as \[
           \lim_{n \to \infty} \frac{a_n}{b_n} = c
       .\] 
       Then, 
       \begin{enumerate}
           \item $c \neq 0$ implies that both $\sum a_n, \sum b_n$ diverge or
               converge.
           \item $c = 0$ and $\sum b_n$ converge, then $\sum a_n$ converge as
               well.
           \item $c = \infty$ and $\sum b_n$ diverge implies that $\sum a_n$ diverge
       \end{enumerate}
   \end{tcolorbox}
\begin{tcolorbox}[title=Ratio Test]
    Let $\sum a_n$ be some series. Then,  \[
        \lim_{n \to \infty} \left| \frac{a_{n+1}}{a_n} \right| = p
    .\] 
    If $p < 1$, the series converges absolutely. If $p > 1$ or infinite, the
    series is divergent. $p = 1$ is inconclusive.
\end{tcolorbox}

\begin{tcolorbox}[title=Root Test]
    Let $\sum a_n$ be some series and  \[
        \lim_{n \to  \infty} \sqrt[n]{|a_n|} = p
    .\] 
    If $p < 1$, the series converges absolutely. If $p > 1$ or infinite, the
    series is divergent. $p = 1$ is inconclusive.
\end{tcolorbox}
\begin{tcolorbox}[title=Alternating Series Test]
    If $a_n \ge 0$ and $a_n$ monotonically decreases to 0, then the alternating
    series  $\sum_{n=1}^\infty (-1)^n a_n$
\end{tcolorbox}
\begin{tcolorbox}[title=Conditional Convergence]
    For any convergent series that is not absolutely convergenct, the series is
    conditionally convergent.
\end{tcolorbox}
\end{multicols}
\end{document}
